% pipeline_steps.tex: the steps of planner.run in Dubins_TSP_New_Implementation_claude, each with what goes
% in and what comes out, in the notation of Dubins_TSP.pdf.
%
% Build from this folder:
%   pdflatex pipeline_steps.tex && pdflatex pipeline_steps.tex
%   mv pipeline_steps.pdf ..
%   rm -f pipeline_steps.aux pipeline_steps.log pipeline_steps.out
\documentclass[10pt]{article}
\usepackage[letterpaper,margin=1in]{geometry}
\usepackage{amsmath,amssymb}
\usepackage{booktabs}
\usepackage{array}
\usepackage{tabularx}
\usepackage{enumitem}
\usepackage{float}
\usepackage[font=small,labelfont=bf]{caption}
\usepackage[hidelinks]{hyperref}

\newcommand{\code}[1]{\texttt{#1}}
\newcommand{\rplan}{r_{\mathrm{plan}}}
\newcommand{\rsafe}{r_{\mathrm{safe}}}
\newcommand{\dsafe}{d_{\mathrm{safe}}}
\newcommand{\dgeo}{d_{\mathrm{geo}}}
\newcommand{\Kmax}{K_{\max}}
\newcommand{\R}{\mathbb{R}}
\newcolumntype{L}{>{\raggedright\arraybackslash}X}
\setlist{nosep,leftmargin=1.5em}
\makeatletter\@secpenalty=0\makeatother   % do not prefer page breaks before headings

% One step: its input, what it does, and its output.
\newcommand{\inputs}{\paragraph{Input.}}
\newcommand{\does}{\paragraph{What it does.}}
\newcommand{\outputs}{\paragraph{Output.}}
\newcommand{\seedone}{\paragraph{On seed 1.}}

\begin{document}
\sloppy

\begin{center}
  {\LARGE The Planner Step by Step: What Goes In and What Comes Out}\\[1.2em]
  {\large November 2026}
\end{center}
\vspace{1em}

{\sloppy
This note follows \code{planner.run(seed, p)} in \code{Dubins\_TSP\_New\_Implementation\_claude/}, the program that
carries out the overall algorithm of Section~0.11 of \code{Dubins\_TSP.pdf}. For every step it states what goes in,
what the step does, and what comes out, with the name and size of each piece of data as the code passes it on. The
steps are numbered 0 to 6, as in the flow diagram \code{output/pipeline\_flow.png}; each step also gives the steps of
Section~0.11 it carries out. The example numbers come from seed~1 with the default parameters (Table~\ref{tab:params}).
\par}

\paragraph{Notation.} The base station is $S_0$ and the sensors are $S_1,\dots,S_n$. In the code, sensor $i$ is counted
from 0, so the paper's $S_{i+1}$ is the code's sensor $i$; in the distance table $D$, row and column $0$ belong to the
base and row and column $i+1$ to sensor $i$. A group of sensors is a set $\mathcal{S}_k$, and a tour $\pi$ is the order
in which one sink visits its group, starting and ending at $S_0$. $B$ is the time budget; in Step~3c the two halves of a
split group are written $A_t$ and $B_t$, as in Section~0.9.

%--------------------------------------------------------------------------------------------------------------------
\section{The Steps at a Glance}

\begin{table}[H]
\centering\small
\renewcommand{\arraystretch}{1.25}
\begin{tabularx}{\textwidth}{@{}l>{\raggedright\arraybackslash}p{3.55cm}>{\raggedright\arraybackslash}p{1.25cm}LL@{}}
\toprule
Step & Code & Sec.~0.11 & Input & Output \\
\midrule
0 & \code{scenario.make} & -- & seed; map settings & map: base, $n$ sensors, dwell times $\tau$, obstacles \\
1 & \code{geodesic.build} & Step 2 & base, sensors, obstacles; margins $\rplan$, $\rsafe+0.5$ &
  distance table $D$ ($(n{+}1)\times(n{+}1)$), predecessors, graph nodes \\
2 & \code{planner.plan\_routes} & Step 1 & all sensor ids & a queue holding one group of all $n$ sensors \\
3a & \code{ga\_tour.optimize} & Steps 3, 6 & one group $\mathcal{S}_k$; positions; $r$; $R$ & visiting order $\pi$ \\
3b & \code{trajectory.build} & Steps 2, 4, 6 & order $\pi$; $D$ and paths; $\tau$; $R$, $v$, $\Delta t$ &
  tour: path, frames, length $L_k$, time $T_k$ \\
3c & \code{partition.bisect} & Steps 5, 7 & $T_k$, $B$; for a split: $\mathcal{S}_k$, $D$, $\tau$ &
  keep the tour as a sink, or two halves $A_t$, $B_t$ to the queue \\
4 & \code{conflict\_game.resolve} & Steps 8--10 & frames, $T_k$, $L_k$ of all $K$ sinks; $\dsafe$, $B$ &
  new frames, $T_k$, $L_k$; one record per game \\
5 & \code{planner.reoptimize} & Step 12 & sinks still in conflict; their tours & new tours; game played again
  (at most 3 rounds) \\
6 & \code{planner.check} & Step 11 & the final result & $K$, $T_{\max}$, $L_{\mathrm{total}}$ and every constraint
  measured \\
\bottomrule
\end{tabularx}
\caption{Every step with its code, its steps in Section~0.11, and what goes in and out. Steps 2--3c repeat until the
queue is empty; Steps 4--5 repeat until no conflict is left or three repairs have been tried.}
\label{tab:glance}
\end{table}

%--------------------------------------------------------------------------------------------------------------------
\section{The Data Passed Between Steps}

\begin{table}[H]
\centering\small
\renewcommand{\arraystretch}{1.25}
\begin{tabularx}{\textwidth}{@{}>{\raggedright\arraybackslash}p{2.4cm}>{\raggedright\arraybackslash}p{2.0cm}L
  >{\raggedright\arraybackslash}p{1.5cm}@{}}
\toprule
Name in the code & Made by & What it holds & Used by \\
\midrule
\code{scen} & Step 0 & \code{base} $\in\R^2$; \code{xy} $\in\R^{n\times2}$, the sensor positions; \code{tau}
  $\in\R^{n}$, the dwell times $\tau_i$; \code{obstacles}, a list of polygons (each an $m\times2$ array of corners)
  & 1, 3, 6 \\
\code{geo} & Step 1 & \code{D} $\in\R^{(n+1)\times(n+1)}$, the geodesic distances $\dgeo$; \code{pred}, the
  predecessors of Dijkstra's runs; \code{nodes} $\in\R^{N\times2}$, the graph's points; \code{obstacles}, the union of
  the original obstacles & 3, 4, 6 \\
$r$ & from $D$ & $r_i = D[0, i{+}1]$, each sensor's distance from the base (eq.~9) & 3a, 3c \\
$D_s$ & from $D$ & $D[1{:}, 1{:}]$, the sensor-to-sensor distances & 3c \\
\code{ids} & Step 2 or 3c & a group $\mathcal{S}_k$: the ids of its sensors & 3a, 3c \\
\code{order} & Step 3a & the tour $\pi$: the group's ids in visiting order & 3b \\
\code{tour} & Step 3b & \code{order}; \code{path}, a dense polyline (points $\Delta s = 0.25$ m apart);
  \code{frames} $\in\R^{F\times2}$, the position every $\Delta t = 0.5$ s from $t=0$; \code{L}, the length $L_k$ [m];
  \code{T}, the mission time $T_k$ [s] & 3c, 4, 5, 6 \\
\code{tours} & Steps 2--3c & the $K$ tours kept, one per sink & 4, 5, 6 \\
\code{result} & Steps 4--5 & \code{tours} (the plan before resolving); \code{frames}, \code{T}, \code{L} after
  resolving; \code{games}; the conflicts \code{before} and \code{after}; the repair \code{round} kept & 6 \\
\code{checks} & Step 6 & $K$, $T_{\max}$, $L_{\mathrm{total}}$, and every constraint measured & the user \\
\bottomrule
\end{tabularx}
\caption{The data that flows through the planner, with its size. $n$ is the number of sensors, $N$ the number of graph
points, $K$ the number of sinks and $F$ the number of frames of a sink.}
\label{tab:data}
\end{table}

\code{planner.run} hands all of it back as one dictionary: \code{scenario}, \code{geo}, \code{result},
\code{checks}, \code{times} and \code{params}.

%--------------------------------------------------------------------------------------------------------------------
\section{Each Step in Detail}

\subsection*{Step 0. Make the map (\code{scenario.make})}

\inputs
\begin{itemize}
  \item The seed, a whole number. The same seed always gives the same map.
  \item From \code{config.py}: a $100\times100$ m square with the base at $(50, 50)$; 10--15 obstacles; 90--110
    sensors; dwell times $\tau_i$ drawn from 1--5 s; 60\,\% of the sensors placed 3--7 m from an obstacle, with half
    of those placement attempts aimed at inner corners.
\end{itemize}
\does
\begin{enumerate}
  \item Places random obstacles (polygons, L, C and cross shapes), each centred at least 25 m from the base, so that
    their $\rplan$-inflated versions neither touch each other nor come within 5 m of the base.
  \item Places the sensors near obstacles, and keeps one only if a sink can drive base~$\to$~sensor~$\to$~base
    keeping $\rsafe$ from every obstacle (this uses Steps 1 and 3b on a small test map).
  \item Spreads the remaining sensors over the free space and draws a dwell time for each.
\end{enumerate}
\outputs
\code{scen}: \code{base}, \code{xy} ($n\times2$), \code{tau} ($n$), \code{obstacles}, and \code{n\_near}, the number of
sensors placed near obstacles. Step~0 is not part of Section~0.11: in the field the map is given.
\seedone 9 obstacles and $n = 98$ sensors.

\subsection*{Step 1. Geodesic distances (\code{geodesic.build}; Section 0.11, Step 2)}

\inputs
\begin{itemize}
  \item \code{base}, \code{xy} and \code{obstacles} from Step 0.
  \item The margins $\rplan = \rsafe + R = 2 + 5 = 7$ m and $\rsafe + 0.5 = 2.5$ m.
\end{itemize}
\does
\begin{enumerate}
  \item Inflates every obstacle by $\rplan$; the corners of the inflated polygons become graph points, after the $n+1$
    terminals ($S_0, \dots, S_n$).
  \item Joins two points by an edge if the straight line between them does not enter an obstacle inflated by
    $\rsafe + 0.5$ (when the line touches a terminal) or by $\rplan$ (between two corners). The edge's weight is its
    length. This is the visibility graph (\code{geodesic.visibility\_graph}).
  \item Runs Dijkstra's algorithm from each of the $n+1$ terminals. Each run gives one row of $D$ and the predecessors
    from which \code{geodesic.path} rebuilds the corner-to-corner path. The graph itself is not kept.
\end{enumerate}
\outputs
\code{geo}: \code{D} ($(n{+}1)\times(n{+}1)$, symmetric, never below the straight-line distance), \code{pred},
\code{nodes}, \code{n\_terminals} $= n+1$, and \code{obstacles}. From $D$ the later steps take $r = D[0,1{:}]$ and
$D_s = D[1{:},1{:}]$. If a sensor cannot be reached, the planner stops with an error.
\seedone 99 terminals, 148 corners, 7898 edges. The longest detour is to $S_{54}$ (the code's sensor 53): $r = 55.0$ m against 41.8 m
in a straight line.

\subsection*{Step 2. The queue of groups (\code{planner.plan\_routes}; Section 0.11, Step 1)}

\inputs The ids $0, \dots, n-1$ of all sensors.
\does Starts a first-in, first-out queue holding one group, $\mathcal{S}_1 = \{0, \dots, n-1\}$: one sink with every
sensor. Steps 3a--3c then take groups from the front of the queue until it is empty.
\outputs The queue $[\mathcal{S}_1]$, an empty list of kept tours, and an empty list of problems.

\subsection*{Step 3a. Visiting order: the genetic algorithm (\code{ga\_tour.optimize}; Section 0.11, Steps 3 and 6)}

\inputs
\begin{itemize}
  \item The group $\mathcal{S}_k$ taken from the front of the queue.
  \item The sensor positions \code{xy}, the base, and $r$ from Step 1.
  \item $R = 5$ m; a population of 50 tours; 100 generations; mutation probability 0.1; the random-number generator
    of Steps 1--7, started from the seed.
\end{itemize}
\does
\begin{enumerate}
  \item Makes 50 start tours that are radially unimodal (Section 0.8.2): each sensor is put at random in $S_{\mathrm{out}}$
    or $S_{\mathrm{in}}$; $S_{\mathrm{out}}$ is visited by increasing $r$, then $S_{\mathrm{in}}$ by decreasing $r$, so
    every start tour has one peak.
  \item Scores each tour by its fitness (Section 0.8.1)
    \[ \mathrm{Fitness}(\pi) = L_{\mathrm{euclid}}(\pi) + m(\pi)\,\pi R, \]
    the straight-line length plus $\pi R \approx 15.7$ m for each radial peak $m(\pi)$. Obstacles enter only through $r$.
  \item In each generation keeps the better 25 tours and adds 25 children: ordered crossover of two kept tours (the
    first $c$ sensors of one, then the rest in the other's order), and with probability 0.1 a swap of two sensors.
    The paper does not specify these two moves; they are the code's choice.
  \item After 100 generations returns the tour with the lowest fitness. A group of one or two sensors needs no search:
    it is visited in order of $r$.
\end{enumerate}
\outputs The visiting order $\pi$ of $\mathcal{S}_k$ (a permutation of its ids). Extra peaks cost $\pi R$ each but are
allowed, so $\pi$ need not be unimodal.
\seedone The first run (all 98 sensors) ends with a best tour of 9 peaks.

\subsection*{Step 3b. The drivable path and its time (\code{trajectory.build}, with \code{dubins\_path}; Section 0.11,
Steps 2, 4 and 6)}

\inputs
\begin{itemize}
  \item The order $\pi$ from Step 3a.
  \item From Step 1: \code{geo} (the corner-to-corner paths through \code{geodesic.path}, and the obstacles).
  \item From Step 0: the dwell times $\tau_i$.
  \item $R = 5$ m, $v = 5$ m/s, $\Delta t = 0.5$ s, $\Delta s = 0.25$ m, $\rsafe = 2$ m, and the Dubins method
    (\code{shortest} by default, or \code{pdf} for the rule of Section 0.13).
\end{itemize}
\does
\begin{enumerate}
  \item \emph{Waypoints:} base, the sensors in order, base, joined by their geodesic paths; the corners of those paths
    become waypoints too. A corner closer than $R$ to a neighbouring waypoint is dropped if the shortcut stays
    $\rsafe + 0.5$ m from every obstacle.
  \item \emph{Headings} (the $\theta$ of Section 0.5): at every waypoint one of 17 candidates (16 directions and the
    bisector of the turn) is chosen by dynamic programming, so that the total Dubins length is as short as possible and
    every leg keeps $\rsafe$ from the obstacles.
  \item \emph{Legs:} every pair of consecutive poses is joined by a Dubins path (\code{dubins\_path.returnDubinsPath}):
    the shortest of LSL, RSR, LSR, RSL, RLR and LRL, or the type of the Section 0.13 rule.
  \item \emph{Time} (eq.~17): driving at constant speed and stopping $\tau_i$ at each sensor,
    \[ T_k = \frac{L_k}{v} + \sum_{i \in \mathcal{S}_k} \tau_i . \]
  \item \emph{Frames:} the position is read every $\Delta t$ seconds from $t = 0$; during a stop the position repeats.
\end{enumerate}
\outputs \code{tour}: \code{order}, \code{path}, \code{frames} ($F\times2$), \code{L} $= L_k$ and \code{T} $= T_k$.
\seedone Sink 8's 5 sensors: 120.5 m of straight lines (what the GA scored), a Dubins path of 176.9 m, and
$T = 176.9 / 5 + 15.6 = 51.0$ s over 103 frames.

\subsection*{Step 3c. The budget, and splitting (\code{partition.bisect}; Section 0.11, Steps 5 and 7)}

\inputs
\begin{itemize}
  \item $T_k$ from Step 3b and the budget $B = 80$ s (eq.~18); $\Kmax = 30$.
  \item For a split: the group $\mathcal{S}_k$, \code{xy}, the base, $D_s$ and $r$ from Step 1, $\tau$, $v$ and $R$.
\end{itemize}
\does
\begin{itemize}
  \item If $T_k \le B$, the tour is kept: one more sink.
  \item If $T_k > B$ and $\mathcal{S}_k$ has more than one sensor, the group is split in two (Section 0.9.1):
    \begin{enumerate}
      \item affinity $W_{ij} = \exp\!\big(-(\dgeo(i,j)/\sigma)^2\big)$, $W_{ii} = 0$, with
        $\sigma = \max(\text{mean nearest-neighbour distance}, R)$;
      \item degrees $d_i = \sum_j W_{ij}$ and the normalised Laplacian $L_{\mathrm{norm}} = I - D^{-1/2} W D^{-1/2}$;
      \item its Fiedler vector $f$, the eigenvector of the second-smallest eigenvalue;
      \item for every threshold $t$, $A_t = \{S_i : f_i \le t\}$ and $B_t = \{S_i : f_i > t\}$; the threshold kept
        minimises $\max\big(\hat T(A_t), \hat T(B_t)\big)$, where $\hat T$ is the length of a sweep round the base in
        polar-angle order with geodesic legs, divided by $v$, plus $\sum\tau_i$.
    \end{enumerate}
    Both halves go to the back of the queue; each later gets its own Steps 3a and 3b. The other groups are not touched.
  \item If a single sensor alone needs $T_k > B$, its tour is kept and the problem is reported.
  \item If the kept tours plus the waiting groups would need more than $\Kmax$ sinks, planning stops and the problem
    is reported.
\end{itemize}
\outputs Either the tour, added to the kept tours, or the two halves $A_t$ and $B_t$, added to the queue. When the queue
is empty, Steps 1--7 hand on \code{tours} ($K$ tours, each with $T_k \le B$) and the list of problems.
\seedone One sink with all 98 sensors would need 1202 s. After 29 GA runs and 14 splits there are $K = 15$ sinks, the
longest taking $T_{\max} = 79.6$ s. The first split gives halves of 52 and 46 sensors.

\subsection*{Step 4. The conflict game (\code{conflict\_game.resolve}; Section 0.11, Steps 8--10)}

\inputs
\begin{itemize}
  \item The \code{frames}, $T_k$ and $L_k$ of all $K$ kept tours (all sinks leave the base at $t = 0$).
  \item \code{geo}, for the obstacle check of a DETOUR.
  \item $\dsafe = 3$ m (eq.~20); the 5 m docking zone round the base; the cost weights $(\beta, \gamma, \delta) =
    (1, 1, 1)$; a longest WAIT of 20 s; a DETOUR side-step of $\dsafe + 0.5 = 3.5$ m; the budget $B$; at most 500
    games.
\end{itemize}
\does
\begin{enumerate}
  \item Finds the earliest moment when two sinks come closer than $\dsafe$ (measured between frames, outside the
    docking zone). The sinks in conflict then are the players.
  \item Gives each player its strategies: GO (keep the plan), WAIT (the shortest stop that clears the conflict, or else the
    least risky one) and DETOUR (step aside and rejoin). A WAIT or DETOUR is offered only if the sink still finishes within $B$.
  \item Each player's cost is
    \[ J_k = \alpha \cdot \text{(seconds closer than $\dsafe$ to any sink)} + \beta\,\Delta T + \gamma\,\Delta L / v +
       \delta\,W, \]
    with $\alpha$ large enough that any danger outweighs any delay. Starting from all GO, the player who gains most
    switches, until no player can gain: a pure Nash equilibrium (Section 0.10.1).
  \item Applies the chosen strategies and moves on to the next conflict.
\end{enumerate}
\outputs New \code{frames}, $T_k$ and $L_k$ (a WAIT adds time; a DETOUR adds time and length), and one record per game:
its players, their strategies, the switches, the potential $\Phi$ after each switch, and the equilibrium. With
\code{--resolver simple}, \code{simple\_resolver.resolve} (Section 0.10.2) takes the same input and gives the same output,
with one record per stop.
\seedone 6 games take the plan from 16 pairs and 63.0 s in conflict to none.

\subsection*{Step 5. Repair (\code{planner.reoptimize}, inside \code{planner.resolve\_and\_check}; Section 0.11,
Step 12)}

\inputs
\begin{itemize}
  \item The pairs still in conflict after Step 4.
  \item The tours of the sinks involved, the map, \code{geo}, $B$, and a second random-number generator, started from
    the seed plus 1000, so that both resolvers start Steps 8--12 alike.
\end{itemize}
\does
\begin{enumerate}
  \item For each sink still in conflict, tries two other tours: its order reversed, and a fresh GA run on its sensors.
    Each is built as in Step 3b, and kept only if it fits $B$ and is less dangerous against the other sinks.
  \item Plays Step 4 again on all the planned tours, with the new ones.
  \item Repeats at most 3 times, stopping as soon as no conflict is left.
\end{enumerate}
\outputs \code{result}, from the round with the least time in conflict: \code{round}, \code{tours} (the plan before
resolving), \code{frames}, \code{T} and \code{L} after resolving, \code{games}, and the conflicts \code{before} and
\code{after}. Conflicts are reported, not guaranteed to be removed: on seeds 1--30 the game leaves some on 8 maps.
\seedone No repair is needed; round 0 is kept.

\subsection*{Step 6. The checks (\code{planner.check}; Section 0.11, Step 11)}

\inputs \code{result} from Step 5, \code{scen}, \code{geo}, the parameters, and the problems found in Step 3c.
\does Measures every objective and constraint of Sections 0.6--0.7 on the final frames, independently of how they were
made.
\outputs \code{checks}:
\begin{itemize}
  \item the objectives: $K$, $T_{\max} = \max_k T_k$, $L_{\mathrm{total}} = \sum_k L_k$ (eqs.~11--14);
  \item \code{coverage\_ok}: every sensor visited exactly once (eqs.~15--16);
  \item \code{budget\_ok}: every $T_k \le B$ (eq.~18);
  \item \code{clearance}: the closest any sink comes to an obstacle, which must be at least $\rsafe$ (eq.~19);
  \item \code{curvature\_x\_R}: the sharpest bend times $R$, which must be at most $1.15$;
  \item \code{conflict\_pairs} and \code{closest}: the pairs still closer than $\dsafe$, and the closest approach
    (eq.~20);
  \item \code{K\_ok} ($K \le \Kmax$) and the list of problems.
\end{itemize}
\seedone $K = 15$, $T_{\max} = 79.6$ s, $L_{\mathrm{total}} = 2927$ m, no conflict left, closest approach 3.16 m.

%--------------------------------------------------------------------------------------------------------------------
\section{How the Steps Connect}

\begin{itemize}
  \item \textbf{Run once:} Steps 0 and 1 at the start, Step 6 at the end.
  \item \textbf{The splitting loop} (Steps 2--3c): take a group from the front of the queue, order it (3a), build its
    path (3b), then keep it or split it (3c). A split puts both halves at the back of the queue. The loop ends when the
    queue is empty, and its output is the $K$ tours.
  \item \textbf{The repair loop} (Steps 4--5): play the game on all $K$ tours (4); if conflicts are left, re-plan the
    sinks involved (5) and play again. At most 3 repairs; the round with the least time in conflict is kept.
  \item \textbf{Three measures of length:} the GA scores straight lines plus a peak penalty (only $r$ is geodesic);
    $\hat T$ in Step 3c uses geodesic legs; only the budget test uses the real Dubins length from Step 3b.
  \item \textbf{Randomness:} one random-number generator, started from the seed, serves Steps 1--7 (Steps 2--3c here);
    a second, started from the seed plus 1000, serves Steps 8--12 (Steps 4--5). The same seed always gives the same
    result.
\end{itemize}

%--------------------------------------------------------------------------------------------------------------------
\section*{Appendix: the Parameters Used}

\begin{table}[H]
\centering\small
\renewcommand{\arraystretch}{1.2}
\begin{tabularx}{\textwidth}{@{}lllL@{}}
\toprule
Symbol & Code & Value & Meaning \\
\midrule
$v$ & \code{v} & 5 m/s & constant speed of a sink \\
$R$ & \code{R} & 5 m & minimum turning radius, wheelbase$/\tan\delta_{\max} = 2.5/0.5$ (eq.~7) \\
$\rsafe$ & \code{r\_safe} & 2 m & obstacle clearance, half the width plus 1 m (eq.~5) \\
$\rplan$ & \code{r\_plan} & 7 m & obstacle inflation for planning, $\rsafe + R$ \\
$\dsafe$ & \code{d\_safe} & 3 m & minimum distance between two sinks (eq.~20) \\
$B$ & \code{B} & 80 s & time budget of every sink (eq.~18) \\
$\Kmax$ & \code{K\_max} & 30 & most sinks allowed \\
-- & \code{base\_radius} & 5 m & docking zone round the base \\
$\Delta t$ & \code{dt} & 0.5 s & time between frames \\
$\Delta s$ & \code{ds} & 0.25 m & spacing of the points of a Dubins path \\
-- & \code{pop\_size}, \code{generations}, \code{mutation} & 50, 100, 0.1 & the GA \\
$(\beta,\gamma,\delta)$ & \code{weights} & $(1, 1, 1)$ & weights of extra time, extra distance and waiting \\
-- & \code{max\_wait}, \code{detour\_margin} & 20 s, 0.5 m & the longest WAIT; DETOUR side-step $\dsafe + 0.5$ \\
-- & \code{curvature\_tol} & 0.15 & allowance when curvature is measured from sampled points \\
-- & \code{max\_games}, \code{repair\_rounds} & 500, 3 & limits of Steps 4 and 5 \\
\bottomrule
\end{tabularx}
\caption{The default parameters, from \code{config.default\_params}.}
\label{tab:params}
\end{table}

\end{document}
